\begin{textsample}{POS Dim 1 – rr – Score 25.00 – rr/rr\_ef\_30.txt}  \label{ex:f1_pos_269}
4.6.2 COMPETÊNCIAS ESPECÍFICAS DE MATEMÁTICA PARA O ENSINO FUNDAMENTAL - Reconhecer \textbf{que} a Matemática é \textbf{uma} ciência humana, fruto das necessidades e preocupações de diferentes culturas, em diferentes momentos históricos, e é \textbf{uma} ciência viva, \textbf{que} contribui para solucionar problemas científicos e tecnológicos e para alicerçar descobertas e construções, inclusive com impactos no mundo do trabalho. - Desenvolver o \textbf{raciocínio} lógico, o \textbf{espírito} de investigação e a capacidade de produzir argumentos convincentes, recorrendo aos conhecimentos matemáticos para compreender e atuar no mundo. - Compreender as relações entre conceitos e procedimentos dos diferentes campos da Matemática ( Aritmética, Álgebra, Geometria, Estatística e Probabilidade ) e de outras áreas do conhecimento, sentindo segurança quanto à própria capacidade de construir e aplicar conhecimentos matemáticos, desenvolvendo a autoestima e a perseverança na busca de soluções. - Fazer observações sistemáticas de aspectos quantitativos e qualitativos presentes nas práticas sociais e culturais, de modo a investigar, organizar, representar e comunicar informações relevantes, para interpretá -las e avaliá -las crítica e eticamente, produzindo argumentos convincentes. - Utilizar processos e ferramentas matemáticas, inclusive tecnologias digitais disponíveis, para modelar e resolver problemas cotidianos, sociais e de outras áreas de conhecimento, validando estratégias e resultados. - Enfrentar situações-problema em múltiplos contextos, incluindo -se situações imaginadas, não diretamente relacionadas com o aspecto prático-utilitário, expressar suas respostas e sintetizar conclusões, utilizando diferentes registros e linguagens ( gráficos, tabelas, esquemas, além de texto escrito na língua materna e outras linguagens para descrever algoritmos, como fluxogramas, e dados ). - Desenvolver e/ou discutir projetos \textbf{que} abordem, sobretudo, questões de urgência social, com base em princípios éticos, democráticos, sustentáveis e solidários, valorizando a diversidade de opiniões de indivíduos e de grupos sociais, sem preconceitos de qualquer natureza. - Interagir com seus pares de forma cooperativa, trabalhando coletivamente no planejamento e desenvolvimento de pesquisas para responder a questionamentos e na busca de soluções para problemas, de modo a identificar aspectos consensuais ou não na discussão de \textbf{uma} determinada questão, respeitando o modo de pensar dos colegas e aprendendo com eles. \textbf{Vivemos} atualmente sob o impacto de profundas transformações na sociedade, provocado pela inserção das tecnologias digitais nos mais diferentes lugares e ambientes.
Isso tem \textbf{exigido} da escola e dos professores \textbf{uma} restruturação do seu papel, ampliado as possibilidades para utilizar esses recursos nos ambientes de ensino.
Além disso, das oito competências específicas da área de conhecimento da matemática, três delas dão ênfase ao uso das tecnologias : reconhecer, utilizar e enfrentar.
No cenário da Educação Matemática, diversos recursos tecnológicos, como computadores, laptops, softwares educativos, calculadoras gráficas, televisão, vídeos, aplicativos da Internet, entre outros, têm favorecido as experimentações matemáticas com tecnologias, dando \textbf{um} novo suporte aos processos de ensino e aprendizagem.
No entanto, conforme afirmam Bonotto e Bisognin ( 2015, p. 2 ), é fato \textbf{que} “ O grande desafio dos professores não é somente utilizar as tecnologias em sala de aula, mais sim, tornar a aula mais envolvente, interativa, criativa e capaz de produzir significados ao aluno por meio da utilização da tecnologia ”.
Desse modo, a integração entre currículo e tecnologia pode conduzir transformações na aprendizagem, nas diretrizes de ensino e nas ações de docentes e \textbf{discentes}, valorizando a formação de conceitos e contribuindo para modificar as práticas de ensino tradicional vigentes ( BORBA ; PENTEADO, 2017 ).
Além disso, atividades realizadas manualmente antes, são executadas rapidamente \textbf{hoje}, por \textbf{intermédio} de recurso tecnológico apropriado para explorar e construir diferentes conceitos matemáticos.
Assim, no aprendizado da Matemática, podem ser explorados diversos recursos tecnológicos \textbf{que} possibilitem aos alunos construir, criar e experimentar diferentes possibilidades, como por exemplo, uso de softwares nas aulas de Geometria e Álgebra, Planilhas Eletrônicas para o estudo e desenvolvimento de probabilidade, estatística e sequências numéricas, além de utilizar os serius games como ferramenta de reforço aos estudos no ensino da matemática, com o uso de sistemas de avaliação por níveis e pontos de experiência, dentre outros.
Para \textbf{que} o aluno alcance as competências propostas na BNCC, o documento propõe \textbf{um} conjunto de habilidades.
Essas habilidades estão relacionadas a diferentes objetos de conhecimento – \textbf{aqui} entendidos como conteúdos, conceitos e processos –, \textbf{que}, por sua vez, são organizados em cinco unidades temáticas correlacionadas, \textbf{que} deverão ser desenvolvidas durante todo o Ensino Fundamental.
Cada \textbf{uma} delas pode receber ênfase diferente, a \textbf{depender} do ano de escolarização. - Números : têm como finalidade desenvolver o pensamento numérico, \textbf{que} \textbf{implica} o conhecimento de maneiras de quantificar atributos de objetos e de julgar e interpretar argumentos \textbf{baseados} em quantidades.
No processo da construção da noção de número, os alunos precisam desenvolver, entre outras, as ideias de aproximação, proporcionalidade, equivalência e ordem, noções fundamentais da Matemática.
Essa unidade temática favorece \textbf{um} estudo interdisciplinar envolvendo as dimensões culturais, sociais, políticas e psicológicas, além da econômica, sobre as questões do consumo, trabalho e dinheiro.
É possível, por exemplo, desenvolver \textbf{um} projeto com a História, visando ao estudo do dinheiro e sua função na sociedade, da relação entre dinheiro e tempo, dos impostos em sociedades diversas, do consumo em diferentes momentos históricos, incluindo estratégias \textbf{atuais} de marketing. - Álgebra : Enfatiza o desenvolvimento de \textbf{uma} linguagem, o estabelecimento de generalizações, a análise da interdependência de grandezas e a resolução de problemas por meio de equações ou inequações.
Nessa perspectiva, é \textbf{imprescindível} \textbf{que} algumas dimensões do trabalho com a álgebra estejam presentes nos processos de ensino e aprendizagem desde o Ensino Fundamental – Anos Iniciais, como as ideias de regularidade, generalização de padrões e propriedades de igualdade.
No entanto, nessa fase, não se propõe o uso de letras para expressar regularidades, por mais \textbf{simples} \textbf{que} sejam.
Já no Ensino Fundamental – Anos Finais, os estudos de Álgebra retomam, aprofundam e ampliam o \textbf{que} foi trabalhado no Ensino Fundamental – Anos Iniciais.
Nessa fase, os alunos devem compreender os diferentes significados das variáveis numéricas em \textbf{uma} expressão, estabelecer \textbf{uma} generalização de \textbf{uma} propriedade, investigar a regularidade de \textbf{uma} sequência numérica, indicar \textbf{um} valor desconhecido em \textbf{uma} sentença algébrica e estabelecer a variação entre duas grandezas.
É necessário, portanto, \textbf{que} os alunos estabeleçam conexões entre variável e função e entre incógnita e equação.
As técnicas de resolução de equações e inequações, inclusive no plano cartesiano, devem ser desenvolvidas como \textbf{uma} maneira de representar e resolver determinados tipos de problema, e não como objetos de estudo em si mesmos. - Geometria : Envolve o estudo de \textbf{um} amplo conjunto de conceitos e procedimentos necessários para resolver problemas do mundo físico de diferentes áreas do conhecimento.
Assim, nessa unidade temática, estudar posição e deslocamentos no espaço, formas e relações entre elementos de figuras planas e espaciais pode desenvolver o pensamento geométrico dos alunos.
Esse pensamento é necessário para investigar propriedades, fazer conjecturas e produzir argumentos geométricos convincentes.
As atividades envolvendo a ideia de coordenadas, já iniciadas no Ensino Fundamental – Anos Iniciais, podem ser ampliadas para o contexto das representações no plano cartesiano, como a representação de sistemas de equações do 1º grau, articulando, para isso, conhecimentos decorrentes da ampliação dos conjuntos numéricos e de suas representações na reta numérica. - Grandezas e medidas : As medidas quantificam grandezas do mundo físico e são fundamentais para a compreensão da realidade.
Assim, essa unidade temática favorece a integração da Matemática a outras áreas de conhecimento, como Ciências ( densidade, grandezas e escalas do Sistema Solar, energia elétrica etc. ) ou Geografia ( coordenadas geográficas, densidade demográfica, escalas de mapas e guias etc. ).
Contribuindo ainda para a consolidação e ampliação da noção de número, a aplicação de noções geométricas e a construção do pensamento algébrico.
Outro ponto a ser destacado refere -se à introdução de medidas de capacidade de armazenamento de computadores como grandeza associada a demandas da sociedade \textbf{moderna}.
Nesse caso, é importante destacar o fato de \textbf{que} os prefixos utilizados para byte ( quilo, mega, giga ) não estão associados ao sistema de numeração decimal, de base 10, pois \textbf{um} quilobyte, por exemplo, corresponde a 1024 bytes, e não a 1000 bytes. - Probabilidade e estatística propõe a \textbf{abordagem} de conceitos, fatos e procedimentos presentes em muitas situações-problema da vida cotidiana, das ciências e da tecnologia.
Assim, todos os cidadãos precisam desenvolver habilidades para coletar, organizar, representar, interpretar e analisar dados em \textbf{uma} variedade de contextos, de maneira a fazer julgamentos bem fundamentados e tomar as decisões adequadas.
Isso inclui raciocinar e utilizar conceitos, representações e índices estatísticos para descrever, explicar e predizer fenômenos.
Merece destaque o uso de tecnologias – como calculadoras, para avaliar e comparar resultados, e planilhas eletrônicas, \textbf{que} ajudam na construção de gráficos e nos cálculos das medidas de tendência \textbf{central}.
No \textbf{que} concerne ao estudo de noções de probabilidade, a finalidade, no Ensino Fundamental – Anos Iniciais, é promover a compreensão de \textbf{que} nem todos os fenômenos são determinísticos.
Para isso, o início da proposta de trabalho com probabilidade está centrado no desenvolvimento da noção de aleatoriedade, de modo \textbf{que} os alunos compreendam \textbf{que} há eventos certos, eventos impossíveis e eventos prováveis.
No Ensino Fundamental – Anos Finais, o estudo deve ser ampliado e aprofundado, por meio de atividades nas quais os alunos façam experimentos aleatórios e simulações para confrontar os resultados obtidos com a probabilidade \textbf{teórica} – probabilidade frequentista.
A progressão dos conhecimentos se faz pelo aprimoramento da capacidade de enumeração dos elementos do espaço amostral, \textbf{que} está associada, também, aos problemas de contagem.
Com relação à estatística, os primeiros passos envolvem o trabalho com a coleta e a organização de dados de \textbf{uma} pesquisa de interesse dos alunos.
No Ensino Fundamental – Anos Finais, a expectativa é \textbf{que} os alunos saibam planejar e construir relatórios de pesquisas estatísticas descritivas, incluindo medidas de tendência \textbf{central} e construção de tabelas e diversos tipos de gráfico.
Esse planejamento inclui a definição de questões relevantes e da população a ser pesquisada, a decisão sobre a necessidade ou não de usar amostra e, quando for o caso, a seleção de seus elementos por meio de \textbf{uma} adequada técnica de amostragem.
Cumpre destacar \textbf{que} os critérios de organização das habilidades na BNCC ( com a explicitação dos objetos de conhecimento aos quais se relacionam e do agrupamento desses objetos em unidades temáticas ) expressam \textbf{um} arranjo possível ( dentre outros ).
Portanto, os agrupamentos propostos não devem ser tomados como modelo obrigatório para o desenho dos currículos.
Essa \textbf{divisão} em unidades temáticas serve tão somente para facilitar a compreensão dos conjuntos de habilidades e de como eles se \textbf{inter-relacionam}.
Nas propostas pedagógicas, devem ser enfatizadas as articulações das habilidades com as de outras áreas do conhecimento, entre as unidades temáticas e no interior de cada \textbf{uma} delas.
Na definição das habilidades, a progressão ano a ano se \textbf{baseia} na compreensão e utilização de novas ferramentas e na complexidade das situações-problema propostas, cuja resolução \textbf{exige} a execução de mais etapas ou noções de unidades temáticas distintas.
Os problemas de contagem, por exemplo, devem, inicialmente, estar restritos àqueles cujas soluções podem ser obtidas pela descrição de todos os casos possíveis, mediante a utilização de esquemas ou diagramas, e, posteriormente, àqueles cuja resolução \textbf{depende} da aplicação dos princípios multiplicativo e aditivo e do princípio da casa dos pombos.
Outro exemplo é o da resolução de problemas envolvendo as operações fundamentais, utilizando ou não a linguagem algébrica.
Levando em consideração \textbf{que} existe \textbf{um} quadro de progressão dessas habilidades ao longo da trajetória escolar, faz -se necessário propiciar contínuas retomadas e constantes avanços no aprendizado matemático, respeitando e estabelecendo conexões entre os anos/séries, entre os segmentos e ao longo de todo o currículo da Educação Básica.
Considerando \textbf{que} os saberes matemáticos \textbf{permeiam} a vida dos estudantes, dentro e fora da escola, faz -se necessário \textbf{que} a comunidade escolar considere os conhecimentos prévios, \textbf{sistematizando} -os de maneira \textbf{dialógica}, ampliando e/ou introduzindo novos conceitos.
Para \textbf{que} isso ocorra, as situações didáticas precisam apresentar desafios, estabelecer conexões entre as unidades da própria Matemática e aproximá -la das outras áreas do conhecimento.
Os estudantes, por sua vez, devem ser estimulados a elaborar hipóteses, testar sua validade, confirmá -las ou refutá -las, discutir estratégias, reelaborar o pensamento, comunicar suas conjecturas, desenvolvendo \textbf{um} \textbf{raciocínio} próprio da atividade matemática e compreendendo como os conceitos se relacionam.

% matched lemmas: abordagem, aqui, atual, basear, central, depender, dialógico, discente, divisão, espírito, exigir, hoje, implicar, imprescindível, inter-relacionar, intermédio, moderno, permear, que, raciocínio, simples, sistematizar, teórico, um, viver
\end{textsample}
